\documentclass[11pt,a4paper]{article}

\usepackage[utf8]{inputenc}
\usepackage[spanish]{babel}
\usepackage[T1]{fontenc}
\usepackage{lmodern}
\usepackage[left=2.5cm,right=2.5cm,top=3cm,bottom=3cm]{geometry}
\usepackage{graphicx}
\usepackage{xcolor}
\usepackage{fancyhdr}
\usepackage{titlesec}
\usepackage{hyperref}

% Configuración de colores
\definecolor{primarycolor}{RGB}{0, 150, 136} % Verde oscuro/Teal médico asociado 
\definecolor{secondarycolor}{RGB}{43, 43, 43}

% Configuración de hipervínculos
\hypersetup{
    colorlinks=true,
    linkcolor=primarycolor,
    urlcolor=primarycolor,
    pdftitle={Manual de Usuario: Médico Asociado},
    pdfauthor={RadiographXpress}
}

\pagestyle{fancy}
\fancyhf{}
\fancyhead[L]{\textcolor{secondarycolor}{\small \textbf{RadiographXpress}}}
\fancyhead[R]{\textcolor{secondarycolor}{\small Médicos Referentes}}
\fancyfoot[C]{\thepage}
\renewcommand{\headrulewidth}{0.4pt}
\renewcommand{\footrulewidth}{0.4pt}

\titleformat{\section}{\Large\bfseries\color{primarycolor}}{\thesection}{1em}{}
\titleformat{\subsection}{\large\bfseries\color{secondarycolor}}{\thesubsection}{1em}{}

\graphicspath{{screenshots/}}

\title{
    \vspace{-2cm}
    \includegraphics[width=0.4\textwidth]{logo_placeholder.png} \\ % Reemplazar con el logo real si existe en Overleaf
    \vspace{1cm}
    \Huge \textbf{RadiographXpress} \\
    \Large \vspace{0.5cm} Manual de Usuario: Médicos Asociados / Referentes
}
\author{Clínica RadiographXpress}
\date{\today}

\begin{document}

\maketitle
\thispagestyle{empty}

\begin{abstract}
\noindent El módulo de Médico Asociado está diseñado para acelerar y optimizar la comunicación inter-institucional. Este portal permite recibir y consultar de forma segura y en tiempo real los resultados de imagenología de tus pacientes directamente dictaminados por los radiólogos de nuestra sucursal.
\end{abstract}

\newpage
\tableofcontents
\newpage

\section{Registro, Validación e Ingreso}

\subsection{Creación de Cuenta}
\begin{enumerate}
    \item Desde la página inicial en la web, selecciona la opción \textbf{"Soy Médico Asociado"} y finaliza en \textit{Crear cuenta de Médico}.
    \item Completa los detalles de tu perfil profesional de forma concisa. Se te solicitará estrictamente ingresar de forma correcta tu \textbf{Cédula Profesional} y \textbf{Especialidad}.
\end{enumerate}

\subsection{Validaciones y Aprobación (Normativas Privacidad)}
Por medidas de seguridad sanitaria nacional, un doctor asociado recién creado carece de permisos de visualización hasta cursar filtros administrativos:
\begin{itemize}
    \item \textbf{Verificación Email:} Recibirás un correo electrónico de conformación. Debes certificar que tu email es correcto presionado al enlace adjunto.
    \item \textbf{Activación Manual Médica:} El personal administrativo central cruzará tu Cédula con el registro gubernamental previo a entregarte tu licencia final. Recibirás otra notificación cuando tu cuenta sea funcional y avalada.
\end{itemize}

\section{El Dashboard Principal}
Superados los filtros, accederás al panel donde se agrupa tu flujo de actividad constante.

\begin{figure}[h]
    \centering
    \includegraphics[width=0.9\textwidth]{associate_doctor_dashboard.png}
    \caption{El Panel (Worklist) del Médico Asociado.}
\end{figure}

\subsection{Consentimiento del Paciente}
\textbf{Importante.} Como Médico Asociado, \textit{no investigas en una base de datos ni abres pacientes libremente}. Por motivos éticos y por secreto profesional, \textbf{el paciente debe buscarte en su plataforma y concederte acceso explícito}. 
Hasta que no seas habilitado por el paciente propio, la lista de tu dashboard permanecerá sin entidades.

\section{Visualización y Revisión Clínica}

Al concederte el acceso un paciente, ocurrirán múltiples transferencias automatizadas del sistema a tu dashboard principal.

\subsection{Estados e Insignias (In-Progress)}
Visualizarás los cuadros de estudio en proceso:
\begin{enumerate}
    \item \textbf{En Progreso:} Indicativo con logo bloqueado. El técnico completó la captura del estudio y un radiólogo lo tiene actualmente activo dictaminando en nuestra base. Sin resultados parciales.
    \item \textbf{Completado:} Estudio con licencia firme lista para inspección.
\end{enumerate}

\subsection{Expediente y Dictamen Clínico Integral}
Abre un estudio completado y analiza en conjunto el veredicto de tu homólogo especialista:
\begin{itemize}
    \item \textbf{Descripción Clínica / Motivo de estudio.}
    \item \textbf{Hallazgos visualizados.}
    \item \textbf{Conclusiones diagnósticas detalladas.}
    \item \textbf{Recomendaciones u observaciones adicionales.}
\end{itemize}

\subsection{Extracción Certificada (PDF)}
En la vista final, puedes oprimir el botón de \textbf{Descargar PDF}. RadiographXpress generará un dictamen oficial membretado y sellado criptográficamente con la firma autógrafa digital de nuestro doctor en turno para adjuntar a tu propio expediente u hospitales foráneos.

\vspace{1cm}
\begin{center}
\textit{Revocaciones: Considera que el paciente puede y tiene la facultad de oprimir "Desvincular" en cualquier momento. De ocurrir así, todos los historiales y archivos visuales desaparecerán de forma hermética de tu cuenta al refrescar tu pantalla.}
\end{center}

\end{document}
